\section{Alternative cancellation estimates}
\label{molb:sec:variants}

The first two comparisons use different information about the same
crossing coefficients. Fractional matching assigns cancellation to the
slack of individual outcomes. A cross-law comparison instead keeps the
smaller absolute mass supplied by the two endpoints in each cone.
We then record the losses from using only one ordered sign conflict
and from retaining weaker projection and trimming estimates.
These are comparisons of proof mechanisms and error budgets; their
lens bounds need not improve the main theorem.

The matching and cross-law arguments below concern the finite-height
forest space $X=J^*$. The ordered and weaker-trimming arguments concern
$B_\infty$, with their own hypotheses stated at the point of use.
Write $\mathbf1_I=d_I$ for a segment functional. The atomic, gauge,
and coordinate-polar descriptions of $X$ agree by
Proposition~\ref{presentation:identity}. Its finite head normers and
atomic approximations come from Lemma~\ref{lem:finite-hull} and
Corollary~\ref{cor:closed-atomic-ball}.

For a finite ancestor-closed head \(H\), put \(P=P_H\) and \(Q=I-P\).
The minimal excluded nodes \(g\) are its gates, and \(T_g\) denotes the
full descendant cone at \(g\).  The contractive projections and the
cone formula of Proposition~\ref{presentation:subtrees} imply the two
estimates used below: gate-coordinate restriction is contractive into
\(\ell_2\), and a combination of rooted tail segments with absolute
coefficient totals \(a_g\) has norm at most \(\|(a_g)\|_2\).
The latter follows by applying the triangle inequality in each cone
and then the cone formula.  Countable cone sums and norm limits use
the completed version of that proposition.

\subsection{Fractional matching of opposite masses}
\label{molb:sec:matching}

\begin{lemma}\label{molb:lem:matching}
Fix a primal unit vector \(z\) supported on \(H\), and a finite
probability distribution of represented atoms
\(A_\omega=\sum_Ic_{\omega I}\mathbf1_I\), with probabilities
\(p_\omega\).  Write
\[
 b_{\omega I}=z(I),\qquad
 L_\omega=1-\sum_Ib_{\omega I}^2,\qquad
 B_\omega=\sum_Ib_{\omega I}Q\mathbf1_I.
\]
At each gate let \(P_g\) and \(M_g\) be the positive and absolute
negative coefficient masses in \(\mathbb E B\).  Then
\begin{equation}\label{molb:eq:matching-capacity}
 \sum_g\min(P_g,M_g)^2\le\mathbb E L.
\end{equation}
\end{lemma}

\begin{proof}
Only crossing segments contribute to \(B_\omega\).  At any gate an
outcome supplies at most one crossing, with signed coefficient
\(b_{\omega g}\).  Regard its absolute mass
\(p_\omega|b_{\omega g}|\) as a supply.  Match positive and negative
supplies at this gate by repeatedly assigning the smaller of the two
remaining supplies to a pair and deleting every exhausted supply.  This
finite procedure produces pair masses \(m\ge0\) whose total is
\(\min(P_g,M_g)\).  The total mass using any one outcome at this gate
is at most its supply.

In a pair of positive mass, the two head intervals end at the parent of
the gate.  They are strictly nested: equality would give equal primal
evaluations, whereas the signs are opposite.  Denote the absolute
evaluations of the longer and shorter heads by \(u>0\) and \(v>0\),
respectively.  Splitting the longer crossing immediately before the
shorter head gives a squared-test gain
\[
 D=(u+v)^2+v^2-u^2\ge uv.
\]
Charge \((m/u)D\) to the outcome containing the longer crossing.  This
charge is at least \(mv\).

Here is the capacity argument for all these charges.  Fix an outcome
\(\omega\) and one of its crossing segments \(I\).  Its absolute
evaluation \(u\) is fixed.  The sum of the fractions \(m/u\) over pairs
in which this segment supplies the longer head is at most \(p_\omega\).
Consequently its total charge is at most \(p_\omega\) times the largest
proposed gain on this segment, or zero if there is no proposal.  Choose
one split giving this largest gain on every affected segment of the
outcome.  These segments were disjoint, so all chosen splits can be made
together.  Their total gain is at most \(L_\omega\), because the new
family is a legal test on the primal unit vector.  Summing gives a bound
\(p_\omega L_\omega\) for the charge to this outcome.  Thus
\[
 \sum_{g}\sum_{\text{pairs at }g}mv\le\mathbb E L.
\]

At a fixed gate the shorter member of a pair contributes the quantity
\(m/v\).  The supply constraint for each outcome implies that the sum
of these quantities over pairs having that outcome as shorter member is
at most its probability.  There is at most one crossing per outcome at
the gate, hence \(\sum m/v\le\sum_\omega p_\omega=1\).
Cauchy--Schwarz now gives
\[
 \min(P_g,M_g)^2=(\sum m)^2
       \le(\sum mv)(\sum m/v)\le\sum mv.
\]
Sum over gates to obtain \eqref{molb:eq:matching-capacity}.
\end{proof}

The independent-sample argument gives a stronger numerical bound for
these same masses. Indeed Lemma~\ref{stab:coef:head-conflict} gives
$4\sum_gP_gM_g\le2\mathbb EL$, so
\[
 \sum_g\min(P_g,M_g)^2\le\sum_gP_gM_g\le\tfrac12\mathbb EL.
\]
The fractional proof above instead displays a matching and charges
its pairs to individual segment-test slacks. The next application
retains its stated capacity bound to show how that argument controls
a whole lens.

\begin{proposition}\label{molb:prop:matching-lens}
For \(0<h<1/4\) and every unit \(x\in X\), there is a finite
ancestor-closed \(H\), depending only on \(x,h\), such that
\(\|x\pm y\|\le1+h\) implies
\begin{equation}\label{molb:eq:matching-lens}
 \|Qy\|\le h+(1+h)
  \bigl[h+2\sqrt{6h}+2(2h+3\sqrt{6h})\bigr].
\end{equation}
\end{proposition}

\begin{proof}
Choose a finitely supported primal unit vector \(z\) with
\(x(z)>1-h\), and enlarge its ancestor closure so that \(\|Qx\|<h\).
Approximate each \((x\pm y)/(1+h)\) within \(h\) by a finite mean of
atoms.  Let \(\mathbb E\) be their equal mixture and
\(\mathbb E_\pm\) the two conditional means.  The identity
\[
 2(1-A(z))=(1-\sum c_I^2)+(1-\sum b_I^2)+\sum(c_I-b_I)^2
\]
has three nonnegative terms on the right.  Also
\(\mathbb EA(z)\ge(1-h)/(1+h)-h\ge1-3h\).  Therefore
\[
 \mathbb EL\le6h,\qquad
 \mathbb E\sum(c_I-b_I)^2\le6h,\qquad
 \mathbb E\|QA-B\|\le\sqrt{6h}.
\]
The corresponding conditional replacement error is at most
\(2\sqrt{6h}\).  Since \(\mathbb EA\) approximates \(x/(1+h)\)
within \(h\), gate-coordinate contraction gives
\[
 \|P-M\|_2\le\|\mathbb EB\|\le2h+\sqrt{6h}.
\]
By Lemma~\ref{molb:lem:matching} and the coordinatewise identity
\(P+M=|P-M|+2\min(P,M)\),
\(\|P+M\|_2\le2h+3\sqrt{6h}\).
Each conditional absolute coefficient mass is at most twice the mixed
mass, whence
\(\|\mathbb E_+B\|\le2(2h+3\sqrt{6h})\).
Approximation and replacement bound the tail of \((x+y)/(1+h)\) by
the bracket in \eqref{molb:eq:matching-lens}; subtracting \(Qx\)
proves the result.  The displayed bound tends to zero with \(h\)
uniformly in the unit center.
\end{proof}

\subsection{Comparing the two endpoint laws}
\label{molx:separate-laws}

The tail can also be recovered by comparing the two endpoint laws
separately. This argument keeps the smaller absolute mass in each
subtree. It gives a different proof, with the following constant.

\begin{proposition}\label{molx:minimum-tail}
Let $X=J^*$ be the finite-height forest space defined by
\eqref{eq:segment-norm}, with the atomic and polar norm identifications
of Proposition~\ref{presentation:identity}. For a finite ancestor-closed
set $A\subset\mathcal F$, let $P_A$ be coordinate restriction on $X$
and let $Q_A=I-P_A$.
If
\[
 P_Ax=x,\qquad Q_Ay=y,\qquad \|x+y\|,\|x-y\|\le1,
\]
then, with $d=1-\|x\|$,
\begin{equation}\label{molx:minimum-conclusion}
 \|y\|\le(\sqrt6+2\sqrt2)\sqrt d\le6\sqrt d.
\end{equation}
\end{proposition}

\begin{proof}
Let $R$ be the minimal vertices of $\mathcal F\setminus A$, and
let $U_r$ be the descendant subtree rooted at $r\in R$.
Lemma~\ref{lem:forest-structure} supplies the contractive projections,
and Lemma~\ref{lem:cones} supplies both the exact subtree norm identity
\eqref{eq:cone-decomposition} and the contractive root-coordinate map.

If $x=0$, the conclusion follows from $\|y\|\le1$. Otherwise,
Lemma~\ref{lem:finite-hull} gives a head-supported test $a\in J$ with
$\rho(a)=1$ and $x(a)=\|x\|$. Since $P_Ay=0$, both endpoints
$z^\pm=x\pm y$ have value $1-d$ on $a$.

Write each represented atom as
\[
 m=\sum_i c_i d_{S_i},\qquad \sum_i c_i^2\le1,
\]
with its disjoint segment family included in the data.
Corollary~\ref{cor:closed-atomic-ball} supplies the norm approximations
used below.

Fix $\eta>0$. Choose two finite laws on represented atoms, with
expectations $E_+,E_-$, such that
$\|E_\pm m-z^\pm\|\le\eta$. For each atom put
\[
 b_i=a(S_i),\qquad
 m_a=\sum_i b_i d_{S_i},\qquad
 D(m)=1-\sum_i b_i^2.
\]
The coefficient identity gives
\[
 D(m)+\sum_i(c_i-b_i)^2
 =1+\sum_i c_i^2-2m(a)\le2(1-m(a)).
\]
Since each endpoint has the same evaluation on $a$, each law
separately satisfies
\begin{equation}\label{molx:separate-deficits}
 E_\pm D(m)\le2(d+\eta),\qquad
 E_\pm\sum_i(c_i-b_i)^2\le2(d+\eta).
\end{equation}
Consequently, for $e=\eta+\sqrt{2(d+\eta)}$,
\begin{equation}\label{molx:separate-error}
 \|z^\pm-E_\pm m_a\|\le e.
\end{equation}
The estimates for separate laws here use $y(a)=0$; they do not follow
from a bound on their mixture alone.

Only segments crossing from $A$ into its complement contribute to
$Q_Am_a$. In each $U_r$ there is at most one such segment in a
represented atom, because it must pass through $r$. Let $b_r(m)$ be
its coefficient, or zero if it is absent. The restriction of $m_a$ to
$J(U_r)$ has norm at most $|b_r(m)|$ and value $b_r(m)$ on the unit
vector at $r$. Segments wholly outside $A$, including those in
components disjoint from $A$, have corrected coefficient zero; their
original coefficients have already been included in
\eqref{molx:separate-deficits}.

The truncation comparison in Lemma~\ref{stab:coef:head-conflict},
applied to the test $a$ and the two represented families, gives
\begin{equation}\label{molx:cross-conflict}
 2\sum_{r:\,b_r(m)b_r(m')<0}|b_r(m)b_r(m')|
       \le D(m)+D(m').
\end{equation}
The important choice here is the law of each sample: we compare one
sample from each endpoint, instead of two samples from their mixture.

Now take $m$ from the plus law and $m'$ independently from the minus
law. Set
\[
 L_r^\pm=E_\pm|b_r(m)|,\qquad S_r^\pm=E_\pm b_r(m).
\]
All these arrays have finite support. Independence and
\eqref{molx:separate-deficits}--\eqref{molx:cross-conflict} yield
\begin{equation}\label{molx:cross-law-product}
 \sum_r(L_r^+L_r^- -S_r^+S_r^-)\le4(d+\eta).
\end{equation}
The endpoint sum $z^++z^-=2x$ vanishes outside $A$.
By \eqref{molx:separate-error} and the contractive root-coordinate map,
$\|S^++S^-\|_2\le2e$. Therefore
\[
 \sum_rS_r^+S_r^-=\tfrac14
  \bigl(\|S^++S^-\|_2^2-\|S^+-S^-\|_2^2\bigr)\le e^2,
\]
and hence
\begin{equation}\label{molx:min-mass}
 \sum_r\min(L_r^+,L_r^-)^2
 \le\sum_rL_r^+L_r^-\le4(d+\eta)+e^2.
\end{equation}

Let $y_r=y|_{J(U_r)}$ and let $e_r^\pm$ be the norm of the
restriction of $z^\pm-E_\pm m_a$ to that subtree. Equation~\eqref{eq:cone-decomposition} and
\eqref{molx:separate-error} give $\|(e_r^\pm)_r\|_2\le e$.
The restrictions of $z^\pm$ are $\pm y_r$, so
\[
 \|y_r\|\le\min(L_r^+,L_r^-)+e_r^++e_r^-.
\]
Combining \eqref{eq:cone-decomposition}, the Euclidean triangle inequality, and
\eqref{molx:min-mass}, we obtain the complete error estimate
\begin{equation}\label{molx:minimum-budget}
 \|y\|\le\sqrt{4(d+\eta)+e^2}+2e,
 \qquad e=\eta+\sqrt{2(d+\eta)}.
\end{equation}
Let $\eta$ decrease to zero. This gives
\eqref{molx:minimum-conclusion}, since $\sqrt6+2\sqrt2<6$.
\end{proof}
\subsection{An ordered sign comparison}
\label{molc:oriented-section}

An ordered sign comparison gives a tail estimate in terms of the
head deficit. We work in the infinite-height space $B_\infty$ of
Remark~\ref{presentation:infinite-height}, on $T=T_\infty$, and write
$J(a)=\rho(a)$ for its segment test norm. Finite coefficient arrays are
identified with coordinate functionals. We use the contractive projections
and norming head tests from Lemma~\ref{binf:projections}, the exact finite
atomic representations from Lemma~\ref{binf:finite-hull}, and the exterior
cone estimates from Lemma~\ref{binf:cones}.

\begin{proposition}\label{molc:oriented-lens}
Let $A$ be a nonempty finite ancestral set and let $x,y\in B_\infty$
with $x=P_Ax$ and $P_Ay=0$.
For every $\varepsilon\ge0$,
\[
 \max\{\|x+y\|,\|x-y\|\}\le1,
 \qquad \|x\|\ge1-\varepsilon
 \quad\Longrightarrow\quad
 \|y\|\le (\sqrt2+\sqrt{10})\sqrt\varepsilon.
\]
\end{proposition}

\begin{proof}
First suppose that $y$ is finitely supported and choose a finite ancestral
set $B\supseteq A$ containing its support. Choose a test $a$ supported in
$A$ such that $J(a)=1$ and $\langle x,a\rangle=\|x\|$. Lemma~\ref{binf:projections} supplies such a test when $x\ne0$; when $x=0$, any $J$-unit test
supported in the nonempty set $A$ has the required equality.

Choose a uniform sign $\sigma\in\{-1,1\}$ and, conditionally on that sign,
a finitely supported law of represented molecules
$m=\sum_i\lambda_i\mathbf1_{S_i}$, with $S_i\subset B$ and
$\sum_i\lambda_i^2\le1$, as supplied by Lemma~\ref{binf:finite-hull},
such that
\[
 \mathbb E(m\mid\sigma)=x+\sigma y.
\]
The representation includes the segment family and its coefficients. Thus
$\mathbb Em=x$ and $\mathbb E(\sigma m)=y$. For each outcome define
\[
 u_i=a(S_i),\qquad D=1-\sum_i u_i^2,\qquad
 H=\sum_i(\lambda_i-u_i)^2,\qquad
 \widetilde m=\sum_i u_i\mathbf1_{S_i}.
\]
Both $D$ and $H$ are nonnegative, and
\begin{equation}\label{molc:joint-budget}
 \mathbb E(D+H)
 =\mathbb E\left(1+\sum_i\lambda_i^2-2\langle m,a\rangle\right)
 \le 2(1-\|x\|)\le2\varepsilon.
\end{equation}
The nonempty suffixes $S_i\setminus A$ remain disjoint segments. Therefore,
with $Q=I-P_A$,
\begin{equation}\label{molc:replacement}
 \|Q(m-\widetilde m)\|\le\sqrt H,
 \qquad
 \|y-\mathbb E(\sigma Q\widetilde m)\|
 \le\mathbb E\sqrt H\le\sqrt{2\varepsilon}.
\end{equation}

Let $R$ consist of the minimal nodes of $B\setminus A$. These nodes are
incomparable, and their descendant subtrees intersected with $B$
partition $B\setminus A$.
For each
$r\in R$, let $c_r$ be the coefficient $u_i$ of the segment crossing from
$A$ through $r$, or zero if there is no crossing. A segment entirely
outside $A$ has $u_i=0$, and each segment meets at most one member of $R$.
Consequently $(Q\widetilde m)(r)=c_r$ and
\[
 \sum_{r\in R}|c_r-m(r)|^2\le H.
\]
Since $\mathbb Em=x$ vanishes on $R$, Jensen's inequality gives
\begin{equation}\label{molc:frontier-mean}
 \sum_{r\in R}(\mathbb Ec_r)^2\le\mathbb EH\le2\varepsilon.
\end{equation}

Compare two represented molecules, using primes for the second.
Write $c^+=\max(c,0)$ and $c^-=\max(-c,0)$ coordinatewise.
Retain only conflicts with $c_r>0>c'_r$ in the nonnegative sum of
Lemma~\ref{stab:coef:head-conflict}. Its truncation proof, applied to
these families and the test $a$, gives
\begin{equation}\label{molc:ordered-conflict}
 2\sum_{r\in R}c_r^+(c'_r)^-\le D+D'.
\end{equation}
Thus the ordered estimate uses only one of the two opposite-sign
orientations controlled by the common boundary comparison.
Take two independent copies of the mixture law. Independence and
\eqref{molc:ordered-conflict} imply
\[
 \sum_{r\in R}(\mathbb Ec_r^+)(\mathbb Ec_r^-)
 \le\mathbb ED.
\]
Thus the unsigned means satisfy
\begin{equation}\label{molc:ten-deficit}
 \begin{aligned}
 \sum_{r\in R}(\mathbb E|c_r|)^2
 &=\sum_{r\in R}(\mathbb Ec_r)^2
     +4\sum_{r\in R}(\mathbb Ec_r^+)(\mathbb Ec_r^-)\\
 &\le 2\varepsilon+4\mathbb ED\le10\varepsilon.
 \end{aligned}
\end{equation}
The coefficient $10$ results from retaining just one ordered conflict
and then using the separate bounds on $\mathbb EH$ and $\mathbb ED$
in \eqref{molc:joint-budget}. Using both orientations in the common
boundary comparison gives a smaller unsigned-mass bound; the present
calculation records the cost of the ordered choice.

Only crossing segments contribute to $Q\widetilde m$. For a finitely
supported test $b$, put
\[
 L_r(b)=\max\{|b([r,s])|:s\in B,\ r\preceq s\}.
\]
The segments attaining these finitely many maxima lie in disjoint
subtrees, so $\sum_r L_r(b)^2\le J(b)^2$. Hence
\[
 |\langle\mathbb E(\sigma Q\widetilde m),b\rangle|
 \le\sum_r\mathbb E|c_r|\,L_r(b)
 \le\sqrt{10\varepsilon}\,J(b).
\]
Taking the supremum over $J(b)\le1$ and using
\eqref{molc:replacement} proves the finite-support estimate. If $R$ is
empty, all these tail sums are zero and the same conclusion holds.

For general $y\in B_\infty$, apply the finite estimate to
$x\pm P_By=P_B(x\pm y)$ along the exhaustion in
Lemma~\ref{binf:projections}. These pairs remain in the unit ball, and
$P_By\to y$ in norm, so the same bound passes to the limit.
\end{proof}

\paragraph{The original-center budget.}\label{molc:original-center-budget}
Put $K=\sqrt2+\sqrt{10}$. Given $0<t<1$, choose $0<d<1$ with
$K\sqrt{3d(1+2d)}<t$. For $x\in S_{B_\infty}$, choose a nonempty finite
ancestral set $A$ with $\|x-P_Ax\|\le d$, and choose a norm-one functional
$x^*$ with $x^*(x)=1$. Set $F=\ker P_A\cap\ker x^*$.
If $y\in F$, $\|y\|\ge1$, and both $\|x\pm ty\|\le1+d$, then
\[
 \|P_Ax\pm ty\|\le1+2d,\qquad \|P_Ax\|\ge1-d.
\]
Apply Proposition~\ref{molc:oriented-lens} after dividing both vectors by
$1+2d$, with deficit $3d/(1+2d)$. It gives
\[
 t\|y\|\le K\sqrt{3d(1+2d)}<t,
\]
a contradiction. Both original endpoint norms are at least $1$ because
$x^*(x\pm ty)=1$, so their average exceeds $1+d/2$. This proves the
averaged asymptotic modulus bound $d/2$ with the stated choice of $d$.

\paragraph{The splitting variant and a half-tail budget.}\label{molc:half-tail-budget}
The splitting proof in Lemma~\ref{stab:coef:head-conflict} gives the
same ordered inequality \eqref{molc:ordered-conflict}, since its
squared-test gain is at least twice the absolute product of the two
head sums at each opposite-sign gate.
Using that proof with the same single orientation therefore retains
the bound $10d$ above.
In a finite ancestral set, let $L$ denote the dual norm of $J$, let the head
$x$ have $L(x)=1$, and suppose $L(x\pm z)\le1+\eta$ with $z$ vanishing on
the head. Applying this version with $d=1-1/(1+\eta)$ gives
\begin{equation}\label{molc:finite-ordered-lens}
 L(z)\le(1+\eta)(\sqrt{2d}+\sqrt{10d}).
\end{equation}
For $0<t<1$, choose $\eta>0$ so that the right side is less than $t/2$.
Approximate an arbitrary unit center within $\eta/4$ by a finite unit
center on a head $A$. For any completed tail $y\in\ker P_A$ of norm at
least $1$, choose a finite ancestral $B\supseteq A$ with
$\|P_By\|>1/2$. If both norms at the finite center with perturbation $ty$
were at most $1+\eta$, projection into $B$ and
\eqref{molc:finite-ordered-lens} would give $t\|P_By\|<t/2$.
Both norms at the finite center are at least $1$. Passing back to the
original center therefore gives an average greater than $1+\eta/4$.

\subsection{A weaker trimming budget}
\label{mold:weaker-trimming}
We work in the infinite-height coordinate predual \(B_\infty\), with
its atomic gauge presentation from Lemma~\ref{binf:finite-hull}.
Let \(D\) be a nonempty finite ancestral set, \(x=P_Dx\),
\(\|x\|=1\), and \(z\in\ker P_D\subset B_\infty\).  We use
\(Q_D=I-P_D\). This calculation keeps only the bounds
$\|Q_D\|\le2$ and a squared-test gain bounded below by the absolute product
of the two opposite head sums per conflict, although the earlier
lemmas give $\|Q_D\|\le1$ and at least twice that gain.
We record the resulting larger error budget from these
weaker inputs. Suppose
\[
 0<h<1,\qquad \max\{\|x+z\|,\|x-z\|\}\le1+h.
\]
Then
\begin{equation}\label{mold:weaker-trimming-bound}
 \|z\|\le(1+h)\bigl(2\sqrt{10h}+2\sqrt{2h}\bigr).
\end{equation}
For the proof, first take \(z\) finitely supported and choose a head-supported
norming functional \(f\) for \(x\), using
Lemma~\ref{binf:projections}.  Lemma~\ref{binf:finite-hull} represents
\((x\pm z)/(1+h)\) by finite probability distributions of segment atoms \(m=\sum_i a_i d_{S_i}\).
Put \(b_i=f(d_{S_i})\), \(\widetilde m=\sum_i b_i d_{S_i}\),
\(L=\sum_i(a_i-b_i)^2\), and \(\Delta=1-\sum_i b_i^2\).
Since \(L+\Delta=1+\sum_i a_i^2-2f(m)\) and
\(\mathbb E_\pm f(m)=1/(1+h)\), each conditional distribution satisfies
\[
 \mathbb E_\pm L\le2h,\qquad \mathbb E_\pm\Delta\le2h.
\]
Let \(B_r=\widetilde m(e_r)\) at the roots immediately outside \(D\), and use
\(\mathbb E\) for the equal mixture of the two distributions.  Each segment
meets at most one such root.  Since \(\mathbb E m=x/(1+h)\),
\[
 \sum_r(\mathbb E B_r)^2\le\mathbb E L\le2h.
\]
For independent outcomes, trimming the longer of two opposite-signed heads
preserves disjointness and increases its squared evaluation by at least
\(|B_rB'_r|\).  Retaining this bound on the gain gives
\[
 \sum_{B_rB'_r<0}|B_rB'_r|\le\Delta+\Delta'.
\]
The modifications can be simultaneous: different exterior roots use distinct
segments of either list.  Independence now gives
\[
 \sum_r(\mathbb E|B_r|)^2
 =\sum_r(\mathbb E B_r)^2+
       2\mathbb E\sum_{B_rB'_r<0}|B_rB'_r|
 \le2h+8h=10h.
\]
For either conditional distribution, the corrected tail has norm at most
\(\| (\mathbb E_\pm |B_r|)_r\|_2\le2\sqrt{10h}\).
Indeed, in each exterior component its mean is a positive scalar times a
convex combination of signed segment indicators; products of these finite
combinations produce disjoint segment atoms.  Recovering the uncorrected
conditional tail costs at most
\(\|Q_D\|\mathbb E_\pm\sqrt L\le2\sqrt{2h}\), proving
\eqref{mold:weaker-trimming-bound}.  Approximation in \(\ker P_D\), with an
arbitrarily small increase of \(h\), proves the completed-tail assertion;
the required density is supplied by Lemma~\ref{binf:projections}.

For \(0<t<1\), choose \(h\in(0,1)\) so that the right side of
\eqref{mold:weaker-trimming-bound} is less than \(t\).  At a finite unit head,
every displacement in the tail of norm at least \(t\) makes one endpoint
norm exceed \(1+h\), while both endpoint norms are at least one.  A finite
unit approximation to an arbitrary unit center with error less than
\(h/4\) therefore gives \(\widehat\delta_{B_\infty}(t)\ge h/4\).
